% Chapter 23 -- Claim: Foundation Models Trained on Market Data Beat GARCH
% Requires: amsmath, amssymb, booktabs, array, tabularx, longtable, tcolorbox (breakable, skins), hyperref or url.
% Defines the callout box only if the preamble has not.
\providecolor{fenavy}{HTML}{1B2A6B}
\makeatletter
\@ifundefined{callout}{%
  \newtcolorbox{callout}{breakable, enhanced,
    colback=white, colframe=gray!70, boxrule=0.5pt, arc=1mm,
    borderline west={2.5pt}{0pt}{gray!70},
    left=8pt, right=8pt, top=4pt, bottom=4pt}%
}{}
\makeatother

\providecolor{fenavy}{HTML}{1B2A6B}
\providecolor{feshade}{HTML}{F3F5FA}
\makeatletter
% Shaded box for optional technical material that general readers may skip.
\@ifundefined{technote}{%
  \newtcolorbox{technote}{breakable, enhanced,
    colback=feshade, colframe=fenavy!60, boxrule=0.5pt, arc=1mm,
    left=8pt, right=8pt, top=4pt, bottom=4pt}%
}{}
\makeatother

\chapter{Claim: Foundation Models Trained on Market Data Beat GARCH}
\label{ch:22}

Anyone who has opened a trading app has seen candlestick charts, in which each day's trading is drawn as a small bar showing where the price opened, how high and low it went, and where it closed. A new kind of artificial-intelligence model learns from billions of these bars, much as a chatbot learns from billions of sentences. This chapter asks whether such models forecast market turbulence better than the statistical tools that risk managers have used for decades.

Chapter~\ref{ch:20} tested general-purpose time-series foundation models\index{foundation models}. These models were trained on electricity demand, traffic, weather and many other series, and then used on financial data without further training. One study in that chapter found that training on financial data helped. In 2025 two lines of work followed that lead. The first builds a model from market data alone. Shi et al.\ (2025) introduced Kronos, trained on more than 12 billion price bars from 45 exchanges. They reported that it forecast prices and volatility better than general foundation models, and produced more realistic synthetic market data. The second line keeps a general model but adapts it. Goel, Pasricha, Magris and Kanniainen (2025) repeatedly retrained Google's TimesFM on realized-volatility data. They reported that it beat the econometric models that risk managers have used for decades.

The benchmark both lines must beat is familiar. GARCH and its relatives, tested in Chapter~\ref{ch:07}, forecast volatility well with a handful of parameters, and they have been refined for forty years. The question is whether a model with millions of parameters, trained on market data, can do better. It is also whether the improvement survives a fair test: the same data, the same loss function, and only periods the model could not have seen during training.

The stakes reach risk management and trading. Better volatility forecasts would improve option pricing, Value-at-Risk and volatility targeting. Realistic synthetic data would help test strategies on more history than the market has produced. The claim on trial\index{claims-on-trial method} is: \textbf{foundation models pretrained or fine-tuned on financial market data beat both general-purpose foundation models and econometric benchmarks such as GARCH and HAR in forecasting volatility.}

The lab in Section~\ref{sec:23-lab} replicates the core volatility comparisons of Shi et al.\ (2025) and Goel et al.\ (2025) on free daily market data. It uses Stooq price bars where a free key is available, and a fallback that uses only the Kenneth R. French Data Library. The goal is to check whether models trained on market data really beat GARCH under a fair test.

By the end of the chapter you will be able to:

\begin{itemize}
  \item describe a price bar, and estimate daily volatility from its high, low, open and close;
  \item explain how Kronos turns price bars into tokens, and how a general model is fine-tuned without look-ahead;
  \item judge synthetic market data against the stylized facts of returns;
  \item compare forecasts with the Diebold--Mariano and Giacomini--White tests, using only data after each model's training cutoff.
\end{itemize}


\section{Foundations: market data, fine-tuning and synthetic data}
\label{sec:23-foundations}

\subsection{OHLCV (K-line) bars}

Exchanges summarize trading over each interval, such as a minute, an hour or a day, in a \textbf{price bar}\index{price bar|textbf}. A bar records five numbers. The \textbf{open} is the first trade price in the interval. The \textbf{high} and the \textbf{low} are the highest and lowest prices. The \textbf{close} is the last price. The \textbf{volume} is the number of shares or contracts traded. Together they are called \textbf{OHLCV} data. Some data sets add a sixth number, the \textbf{amount}, which is the value traded in currency. Charts draw each bar as a candlestick. In Chinese and Japanese markets bars are often called \textbf{K-lines}\index{K-line}, which is why Kronos's authors speak of the ``language of K-lines.''

A bar carries more information than a closing price. A day that closes where it opened may still have swung widely between its high and its low, and that swing says something about volatility. The volume says how actively the market traded.

\subsection{Range-based volatility: Parkinson and Garman--Klass}

Chapter~\ref{ch:07} measured a day's variance in two ways. The squared daily return is noisy. Realized variance\index{realized variance} from intraday returns is accurate, but it needs intraday data, which are rarely free. Daily bars allow a middle path. In plain words, how far the price traveled between its high and its low during the day says much more about the day's turbulence than where it happened to close. The shaded box gives the two standard formulas.

\begin{technote}
\textbf{Technical note: range-based variance estimators.} Write $H_t$, $L_t$, $O_t$ and $C_t$ for the high, low, open and close on day $t$. Parkinson (1980) showed that the squared log range estimates the day's variance:
\begin{equation}
  \hat\sigma^2_{P,t} = \frac{1}{4 \ln 2} \left(\ln \frac{H_t}{L_t}\right)^2.
  \label{eq:23-parkinson}
\end{equation}
Garman and Klass (1980) added the open and close:
\begin{equation}
  \hat\sigma^2_{GK,t} = \frac{1}{2} \left(\ln \frac{H_t}{L_t}\right)^2 - \left(2\ln 2 - 1\right) \left(\ln \frac{C_t}{O_t}\right)^2.
  \label{eq:23-gk}
\end{equation}
Under ideal conditions, with continuous trading and no drift, these estimators are several times more precise than the squared close-to-close return. In practice they have two weaknesses. They ignore the jump between yesterday's close and today's open, which matters for stock markets that close overnight. They are also biased downward when trading is not continuous, because the recorded high and low miss the true extremes. They are still a large improvement on squared returns, and this chapter's lab uses the Garman--Klass estimator to measure daily variance.
\end{technote}

\textbf{Worked example.} A stock index opens at 100, trades between 98 and 103, and closes at 101. Then $\ln(H/L) = \ln(103/98) = 0.0498$ and $\ln(C/O) = \ln(1.01) = 0.0100$. The Parkinson estimate is $0.0498^2/(4 \times 0.693) = 0.00089$, a daily volatility of about 3.0\%. The Garman--Klass estimate is $0.5 \times 0.0498^2 - 0.386 \times 0.0100^2 = 0.00120$, a daily volatility of about 3.5\%.

\subsection{Tokenizing multivariate bars}

Chapter~\ref{ch:20} explained how Chronos turns a single series into tokens, by scaling it and sorting each value into one of a fixed set of bins. A price bar has five or six numbers that move together, so Kronos turns the whole bar into one token, a single ``word'' describing the whole day. The shaded box explains how.

\begin{technote}
\textbf{Technical note: how Kronos builds its tokens.} A small encoder network compresses each normalized bar into a short code of zeros and ones, using a method called \textbf{binary spherical quantization}. In the paper's main setting the code has 20 bits, which allows about a million different tokens. The code is split into a \textbf{coarse} part and a \textbf{fine} part of 10 bits each, much as a price can be described first roughly and then precisely.
\end{technote}

A transformer then learns the ``language'' of these tokens. Given a history of bars, it predicts the coarse part of the next token, then the fine part given the coarse part, and so on, one bar at a time. To forecast, Kronos samples many possible future paths in this way and averages them. The same sampling produces synthetic price histories. The paper describes three sizes of the model, with about 25 million, 100 million and 500 million parameters. Each reads a history of up to 512 bars.

\subsection{Fine-tuning: full, incremental and parameter-efficient}

A pretrained model can be adapted to a new task by \textbf{fine-tuning}\index{fine-tuning|textbf}: continuing its training on data from that task. There are three main ways to do it.

\begin{itemize}
  \item \textbf{Full fine-tuning} updates every parameter. It is the most flexible and the most expensive, and with a small data set it can overfit and forget what the model learned in pretraining.
  \item \textbf{Parameter-efficient fine-tuning} freezes most of the network and trains only a small part. One option trains only the last layers. Another, called \textbf{LoRA}\index{LoRA} (low-rank adaptation; Hu et al.\ 2022), adds to each large weight matrix a correction made of two thin matrices, so that only a small number of new parameters are trained.
\item \textbf{Incremental fine-tuning} repeats the training as new data arrive. Goel and coauthors fine-tuned TimesFM on the first half of their sample and tested it on the next fifth. They then added that fifth to the training data, fine-tuned again, tested on the following fifth, and so on. This is the walk-forward design of Section~\ref{sec:09-walkforward}, applied to a neural network.
\end{itemize}

The rule that matters most is the same as in every forecasting chapter: fine-tune only on data that end before the period being forecast. A model fine-tuned on the whole sample and then tested on part of it has seen the answers.

\subsection{Synthetic market data and the stylized facts}

A model that can generate realistic price paths has uses beyond forecasting. Risk managers can use synthetic paths to stress-test portfolios, and researchers can use them to enlarge small data sets. A synthetic path is useful only if it behaves like real data. Cont (2001) summarized the \textbf{stylized facts}\index{stylized facts|textbf} that real returns show across markets, several of which Chapter~\ref{ch:07} documented:

\begin{itemize}
  \item daily returns have fat tails, with far more extreme days than a normal distribution allows;
  \item returns themselves show almost no autocorrelation, but absolute and squared returns show strong, slowly decaying autocorrelation, which is \textbf{volatility clustering}\index{volatility clustering};
  \item volatility tends to rise after prices fall, the \textbf{leverage effect}\index{leverage effect};
  \item returns over longer horizons, such as a month, look closer to normal than daily returns.
\end{itemize}

A GARCH model with fat-tailed shocks reproduces most of these facts with a few parameters, so it is the natural benchmark for any generative model. Kronos's authors measured realism in two ways: how well a classifier could tell synthetic paths from real ones, and how well a model trained only on synthetic data performed on real data.

\subsection{The Giacomini--White test}

The Diebold--Mariano test of Section~\ref{sec:09-accuracy} asks whether one forecast has a lower average loss than another. It assumes that the models' parameters are known, or estimated on ever-growing samples. Models are often re-estimated on a rolling window, or fine-tuned again and again. For that case, Giacomini and White (2006) offer a test that compares whole forecasting methods, estimation included. In plain words, it asks whether one forecasting method is reliably better than another once all the re-estimating is counted, and whether its advantage depends on market conditions.

\begin{technote}
\textbf{Technical note: the Giacomini--White test.} Write $d_{t+1}$ for the difference in losses between methods A and B at time $t+1$. The test asks whether the loss difference can be predicted from information available at time $t$:
\begin{equation}
  H_0: E\left[d_{t+1} \mid \mathcal{F}_t\right] = 0.
  \label{eq:23-gw}
\end{equation}
In practice the researcher chooses a few variables $h_t$ known at time $t$, such as a constant and yesterday's loss difference. A standard Wald test then checks whether the averages of $h_t d_{t+1}$ are all zero. With only a constant, the test reduces to a Diebold--Mariano test, but it stays valid when models are re-estimated on rolling windows. With more variables it can show that one method is better in some conditions, such as calm markets, and worse in others.
\end{technote}

\section{How it works}
\label{sec:23-how}

\textbf{Why train on market data?} A model trained on electricity and weather learns patterns such as daily cycles, trends and smooth seasons. Markets have little of that. Their most reliable patterns are fat tails, volatility clustering, and the link between falling prices and rising volatility. A model trained on billions of price bars sees those patterns constantly and in every market. Kronos's authors argue that this is why it forecasts volatility and generates realistic paths better than general models.

\textbf{Why fine-tuning helps.} Markets change. A model pretrained once, or a general model, may have learned the volatility patterns of a different era or a different domain. Fine-tuning on recent data from the series being forecast lets the model adjust its scale and dynamics. Goel and coauthors found that TimesFM used zero-shot was a reasonable but not leading forecaster of volatility, and that repeated fine-tuning was what made it competitive.

\textbf{A new leakage risk.} A model trained on 45 exchanges has very likely seen the S\&P 500, the major currency pairs and the large stocks that a student would test it on, possibly over the same dates. If so, a test on that history measures memory, not forecasting skill. This is the pretraining leakage\index{leakage!pretraining} that Section~\ref{sec:21-contam} described for general models. It is more severe here, because the training data and the test data come from the same markets. The only clean test uses data that arrived after the model's training data were collected.

\section{The claim}
\label{sec:23-claim}

Stated so that it can fail:

\begin{callout}
\textbf{Foundation models that are pretrained or fine-tuned on financial market data forecast volatility more accurately than GARCH and HAR models, and generate synthetic price paths that match the stylized facts of returns better than a GARCH simulation, on data that arrived after the models were trained.}
\end{callout}

The claim has a \textbf{forecasting form} and a \textbf{generative form}. The tests are whether the comparisons use well-specified \textbf{benchmarks} such as GARCH and HAR on the same data; whether the advantage holds under the QLIKE \textbf{loss} that Chapter~\ref{ch:07} recommended, not only squared error; whether it survives on \textbf{clean data} the model could not have seen; and whether it has \textbf{economic value}, changing a risk number such as a 99\% VaR by enough for its backtest to detect.

\section{The original evidence}
\label{sec:23-evidence}

\subsection{Shi et al. (2025): Kronos}

Shi and coauthors\index{Kronos|textbf} trained Kronos on more than 12 billion bars from 45 exchanges, at seven frequencies with six fields each (open, high, low, close, volume and amount), and compared it with 25 alternatives, from general foundation models and deep networks to GARCH and generative models. Mostly used zero-shot, Kronos had a rank information coefficient (the correlation between predicted and actual rankings of returns across assets) 93\% higher than the best general foundation model and 87\% higher than the best non-pretrained model, a 9\% lower mean absolute error for volatility, a 22\% improvement in synthetic-data fidelity, and good results in a simulated long-only test on Chinese A-shares. These are the authors' own benchmarks, and the main text does not make clear which periods and markets were held out of pretraining. The code and smaller versions are open, so independent testing is possible.

\subsection{Goel, Pasricha, Magris and Kanniainen (2025)}

Goel and coauthors forecast next-day realized variance, from five-minute returns in the Oxford-Man Realized Library, for 21 stock indices from 2000 to 2021. They compared TimesFM 2.0, zero-shot and incrementally fine-tuned, with realized GARCH, ARFIMA, HAR and CHAR models in levels and logs, scored with six losses including MSE and QLIKE, and tested with Diebold--Mariano and Giacomini--White tests and the model confidence set, which identifies the models that cannot be told apart from the best.

Fine-tuned versions generally beat zero-shot ones. A fine-tuned TimesFM reading 64 days of log variance had the lowest average squared error and was never significantly worse than any model under three losses. Under QLIKE, the 512-day TimesFM versions, zero-shot and fine-tuned, beat every econometric model at the 10\% level. Realized GARCH remained hard to beat at predicting whether volatility would rise or fall.

\subsection{Marconi (2025)}

Marconi tested Tiny Time Mixers (TTM) and Chronos on Treasury yield changes, euro--dollar volatility and an equity spread. Fine-tuned TTM beat an identical network trained from scratch by 25\% to 50\% with limited data and 15\% to 30\% with longer data, reaching similar accuracy with three to ten fewer years. But specialized traditional models matched or beat it in two of the three tasks. Pretraining saves data; it does not beat the best specialized model.

\section{Lab~\thechapter: Replicating Kronos and fine-tuned TimesFM on daily market data}
\label{sec:23-lab}

\textbf{Claim on trial\index{claims-on-trial method}:} foundation models trained or fine-tuned on market data forecast volatility better than GARCH and HAR. The lab replicates Shi et al.\ (2025) and Goel, Pasricha, Magris and Kanniainen (2025) with the open small version of Kronos, TimesFM, Chronos-Bolt from Lab~\ref{ch:20}, GARCH(1,1) and HAR. It uses daily OHLCV bars for major stock indices and exchange-traded funds from a free source such as Stooq, daily market returns from the Kenneth R. French Data Library\index{Kenneth French data library} and the VIX from FRED\index{FRED}. Because intraday data are not free, daily variance is measured by the Garman--Klass estimator of Equation~\eqref{eq:23-gk}, and every model forecasts tomorrow's value from data up to today. The zero-shot parts run on a laptop; fine-tuning TimesFM needs a GPU, such as one in a free cloud notebook.

\begin{itemize}
  \item \textbf{Part A, zero-shot forecasts.} Kronos, TimesFM and Chronos-Bolt forecast next-day returns, scored by rank correlation across assets, and next-day variance, scored by MSE and QLIKE.
  \item \textbf{Part B, the Goel et al.\ replication.} TimesFM is fine-tuned on log variance and tested against GARCH and HAR with Diebold--Mariano and Giacomini--White tests.
  \item \textbf{Part C, fine-tuning choices.} History lengths of 64, 128 and 512 days, and fine-tuning of the last layers only or with LoRA.
  \item \textbf{Part D, synthetic paths.} One-year S\&P 500 paths from Kronos and from a Student-$t$ GARCH, scored on kurtosis and volatility clustering.
  \item \textbf{Part E, risk.} Each variance forecast becomes a 99\% one-day VaR, backtested with the Kupiec and Christoffersen tests of Chapter~\ref{ch:08}.
  \item \textbf{Part F, clean data.} Parts A, B and E repeated on days after each model's public release only. This is the test of the claim (Lab Figures~\ref{ch:22}.1 to~\ref{ch:22}.3).
\end{itemize}

\textbf{A fallback that needs no download.}\index{squared return!as variance proxy} If the OHLCV bars cannot be obtained, the lab uses only the daily market return and 49 industry portfolios from the French library. The day's variance becomes the squared close-to-close return, scoring uses QLIKE in the form $\log F + RV/F$ of Section~\ref{sec:10-eval}, which accepts a zero value, and Kronos, which needs full bars, is left out. The squared return is a much noisier measure, so fewer differences will be significant, but with a robust loss it does not on average change which model ranks first (Patton 2011). A clear win in the fallback mode is therefore likely to hold with better data, but a tie is not evidence that the models are equally good.

% [Draft note: the figures and numbers are generated by running the lab on data at the time of publication. The book prints the results; the text above does not anticipate them.]

\label{sec:23-stress}
\textbf{Robustness checks.} Losses are reported separately for February to April 2020 and for 2022, asking whether any model reacts before the jump or only after it. On simulated GARCH data, GARCH is the correct model, so a foundation model that beats it is fitting noise. A large drop in accuracy on small-country indices would suggest that part of the major-index performance is memory. Finally, HAR takes a fraction of a second while a fine-tuned foundation model needs a GPU and hours, so the lab reports the QLIKE gain per hour of computation. The lab ends by asking you to reach your own verdict on the scale of Section~\ref{sec:01-method} and compare it with Section~\ref{sec:23-verdict}.

\section{Where the evidence stands}
\label{sec:23-verdict}

The published studies point to an early picture: \textbf{training on market data appears to help, and a general model fine-tuned on volatility data can match and sometimes beat the best econometric models. But the evidence is young and mostly from working papers, and none of it yet uses data the models could not have seen.} The table summarizes the evidence as of November 2026.

{\small
\begin{longtable}{@{}>{\raggedright\arraybackslash}p{0.19\linewidth}>{\raggedright\arraybackslash}p{0.13\linewidth}>{\raggedright\arraybackslash}p{\dimexpr 0.45\linewidth-6\tabcolsep\relax}>{\raggedright\arraybackslash}p{0.23\linewidth}@{}}
\toprule
Situation & Holds? & Why & Evidence \\
\midrule
\endfirsthead
\multicolumn{4}{@{}l}{\emph{(continued)}} \\
\toprule
Situation & Holds? & Why & Evidence \\
\midrule
\endhead
\midrule
\multicolumn{4}{r@{}}{\emph{(continued on the next page)}} \\
\endfoot
\bottomrule
\endlastfoot
Kronos, volatility & \textbf{Partly} & 9\% lower error on the authors' own benchmarks; not verified independently. & Shi et al.\ (2025) \\
\addlinespace[3pt]
Kronos, synthetic data & \textbf{Partly} & 22\% better fidelity on the authors' own measure; not verified independently. & Shi et al.\ (2025) \\
\addlinespace[3pt]
Fine-tuned TimesFM & \textbf{Yes} & Beat GARCH and HAR under several losses, including QLIKE, on 21 indices. & Goel, Pasricha, Magris and Kanniainen (2025) \\
\addlinespace[3pt]
Zero-shot TimesFM & \textbf{Partly} & Behind the fine-tuned versions, but ahead of the econometric models under QLIKE with a 512-day history. & Goel, Pasricha, Magris and Kanniainen (2025) \\
\addlinespace[3pt]
Direction of volatility & \textbf{No} & Realized GARCH remained hard to beat. & Goel, Pasricha, Magris and Kanniainen (2025) \\
\addlinespace[3pt]
Against specialized models & \textbf{Weakly} & Specialized models matched or beat it in two of three tasks; pretraining mainly saved data. & Marconi (2025) \\
\addlinespace[3pt]
Clean post-release data & \textbf{Untested} & No study yet reports such results. & None yet \\
\end{longtable}
}

\textbf{What would change the picture.} The claim would gain support if foundation models beat GARCH and HAR significantly under QLIKE on post-release data, by enough to change a VaR backtest, and lose it if their advantage disappeared there.

\label{sec:23-open}
\textbf{Open questions.} How much accuracy on major indices comes from having seen them in pretraining, and can synthetic paths from generative models improve stress tests and capital estimates?

\section{What practitioners and investors can use}
\label{sec:23-investors}

Keep HAR and GARCH as the yardstick, because they are cheap, transparent and hard to beat, and admit a foundation model into a risk system only if it beats them out of sample, under QLIKE, by enough to change a risk decision. Adapted models beat models used as delivered, but every fine-tuning step must use only data available at the time. Synthetic paths can widen a stress test beyond history's few crises, but only if they reproduce fat tails and volatility clustering as well as a GARCH simulation. Ask a vendor what data the model was trained on and up to what date, and how it compares with HAR and GARCH after that date.

\begin{callout}
\textbf{What Failed: forecasting volatility from calm data, March 2020.} In mid-February 2020 the VIX was below 15, and models estimated on the calm years before forecast low volatility. Within a month the S\&P 500 fell by about a third, and on March 16 the VIX closed at 82.69, its highest close on record. Every model that forecasts from past data, from GARCH to a foundation model, raised its forecast only after the turmoil began. The claim that a model trained on calm data can anticipate a sudden change of regime \textbf{fails}. A better model can offer only a faster, more accurate response once the change begins, which is what this chapter's crisis tests measure.
\end{callout}
